\chapter{Integrals depending on a parameter}
\SourceRef{Zorich II, Sections 17.1--17.2, PDF pp. 421--442, especially the uniform difference-quotient proof and the base-point version of improper differentiation; Apostol, Section 7.23}
The integration variable is real and the parameter is real. Scalar values may
be real or complex, with all complex derivative estimates understood through
the real-parameter norm inequality. The results here concern proper Riemann
integrals and uniform improper tails; no Lebesgue dominated-convergence or
general Fubini theorem is an implicit prerequisite.

\KEntry{K03.14.proper_definition}{A proper integral family}
\begin{definition}
Let $T\subset\R$ and let $f:[a,b]\times T\to\K$. If each section
$x\mapsto f(x,t)$ is Riemann integrable, define
$F(t)=\int_a^bf(x,t)\dd x$. The formula defines $F$ only on the
parameters where that section is integrable; the symbol alone supplies
no existence or parameter regularity.
\end{definition}

\KEntry{K03.14.proper_continuity}{Continuity under a finite integral}
\begin{theorem}
Suppose the sections are integrable near $t_0$ and
$\sup_{x\in[a,b]}|f(x,t)-f(x,t_0)|\to0$ as $t\to t_0$ in $T$.
Then $F$ is continuous at $t_0$. In particular joint continuity on a
compact rectangle $[a,b]\times[c,d]$ gives a continuous $F$ on $[c,d]$.
\end{theorem}
\begin{proof}
The finite norm estimate gives
\[|F(t)-F(t_0)|\le(b-a)\sup_x|f(x,t)-f(x,t_0)|\longrightarrow0.
\]
For joint continuity, compact uniform continuity gives the stated
uniform-in-$x$ estimate whenever $|t-t_0|$ is sufficiently small.
The same compactness gives continuous, hence integrable, $x$-sections.
At parameter endpoints continuity is relative to the parameter interval.
\end{proof}

\KEntry{K03.14.proper_derivative}{Differentiation under a finite integral}
\begin{theorem}
Let $J$ be an open parameter interval. Suppose $f$ and its real partial
derivative $f_t$ are jointly continuous on $[a,b]\times J$.
Then $F\in C^1(J)$ and
\[F'(t)=\int_a^bf_t(x,t)\dd x.
\]
It is enough to have these conditions on every compact parameter
neighborhood on which the conclusion is used.
\end{theorem}
\begin{proof}
Fix $t_0$ and a compact neighborhood $[c,d]\subset J$.
Uniform continuity of $f_t$ there gives a modulus $\omega(\delta)\to0$
for changes of parameter distance at most $\delta$, uniformly in $x$.
For $h\ne0$ with $t_0+h\in[c,d]$, the norm mean-value inequality
applied to $s\mapsto f(x,s)-s f_t(x,t_0)$ gives
\[\left|\frac{f(x,t_0+h)-f(x,t_0)}h-f_t(x,t_0)\right|
\le\sup_{s\text{ between }t_0,t_0+h}|f_t(x,s)-f_t(x,t_0)|
\le\omega(|h|).
\]
This is valid for complex scalars as a norm inequality, without claiming
a single complex Lagrange equality. Integrate and use the finite norm
bound to obtain difference-quotient error at most $(b-a)\omega(|h|)$.
It tends to zero, proving the derivative formula. The preceding continuity
theorem applied to $f_t$ proves continuity of the resulting derivative.
\end{proof}

\KEntry{K03.14.moving_endpoints}{Leibniz's rule with moving endpoints}
\begin{theorem}
Let $\alpha,\beta\in C^1(J,\R)$. Near a parameter $t_0$, suppose
$f,f_t$ are jointly continuous on a rectangle containing all interval
endpoints $\alpha(t),\beta(t)$ and the intervals between them. Then
\begin{align*}
\frac{\dd}{\dd t}\int_{\alpha(t)}^{\beta(t)}f(x,t)\dd x
={}&f(\beta(t),t)\beta'(t)-f(\alpha(t),t)\alpha'(t)\\
&+\int_{\alpha(t)}^{\beta(t)}f_t(x,t)\dd x.
\end{align*}
Orientation permits the two endpoints to cross.
\end{theorem}
\begin{proof}
Choose a fixed regular $x_*$ in the spatial interval and set
$H(y,t)=\int_{x_*}^yf(x,t)\dd x$. The first fundamental theorem gives
$H_y(y,t)=f(y,t)$, and the preceding differentiation theorem gives
$H_t(y,t)=\int_{x_*}^yf_t(x,t)\dd x$ for fixed $y$.
For completeness along a moving endpoint $y(t)$, split
\begin{align*}
H(y(t+h),t+h)-H(y(t),t)
={}&\int_{y(t)}^{y(t+h)}f(x,t+h)\dd x\\
&+\int_{x_*}^{y(t)}(f(x,t+h)-f(x,t))\dd x.
\end{align*}
The first quotient tends to $f(y(t),t)y'(t)$: subtract that constant
integrand and bound the error by
$|y(t+h)-y(t)|/|h|$ times the local uniform continuity error.
The second quotient tends to $\int_{x_*}^{y(t)}f_t$ by the uniform
difference-quotient estimate. Subtract these formulas for $y=\beta$
and $y=\alpha$ and use oriented additivity. This establishes the rule
using one-dimensional estimates only.
\end{proof}

\KEntry{K03.14.proper_order}{Interchanging two finite integrations}
\begin{theorem}
For continuous scalar $f$ on $[a,b]\times[c,d]$ both iterated proper
integrals exist and
\[\int_c^d\left(\int_a^bf(x,t)\dd x\right)\dd t
=\int_a^b\left(\int_c^df(x,t)\dd t\right)\dd x.
\]
\end{theorem}
\begin{proof}
The inner sections are integrable and their integral functions are
continuous by the continuity theorem, so the outer integrals exist.
Let $\omega(\delta)$ be a uniform continuity modulus on the rectangle.
Take tagged partitions $P,Q$ in $x,t$ with meshes below $\delta$ and
form the finite double sum
$D=\sum_{i,j}f(\xi_i,\tau_j)\Delta x_i\Delta t_j$.
For every fixed $t$, replacing $\int_a^bf(x,t)\dd x$ by
$\sum_i f(\xi_i,t)\Delta x_i$ has error at most
$(b-a)\omega(\delta)$, uniformly in $t$, by integrating the
cellwise difference. Integrating that error in $t$, and then replacing
each $\int_c^df(\xi_i,t)\dd t$ by its tagged sum, shows the first
iterated integral differs from $D$ by at most
$2(b-a)(d-c)\omega(\delta)$. Reverse the order to obtain the same
bound for the other iterated integral. Their difference tends to zero
as $\delta\downarrow0$. This is an equality of limits of finite sums,
not an application of a general multiple-integration theorem.
\end{proof}

\KEntry{K03.14.uniform_improper}{Uniform improper convergence}
\begin{definition}
Assume local proper integrability in $x$ for every $t\in T$ and write
$F_R(t)=\int_a^Rf(x,t)\dd x$. The family converges uniformly
improperly on $T$ if there is $F:T\to\K$ with
$\sup_{t\in T}|F_R(t)-F(t)|\to0$ as $R\to\infty$.
Equivalently, for every $\eps>0$ there is $R_0$ such that
\[\left|\int_R^Sf(x,t)\dd x\right|<\eps
\quad(t\in T,\ S\ge R\ge R_0).
\]
For a finite right singular endpoint $b$, require the same condition
for $r_0<R<S<b$ with $r_0$ sufficiently close to $b$; for a left
endpoint reverse the truncations. Several singularities require
independent uniform conditions at every endpoint.
\end{definition}
\begin{proof}[Equivalence]
Uniform convergence implies the displayed estimate by the triangle
inequality. Conversely, for each fixed $t$ the ordinary Cauchy criterion
gives a finite limit $F(t)$. Hold $R\ge R_0$ and let $S\to\infty$
in the displayed inequality. This gives $|F(t)-F_R(t)|\le\eps$
simultaneously for all $t\in T$. Apply the condition with an arbitrarily
smaller tolerance to obtain uniform convergence. The same argument works
at a finite singular endpoint.
\end{proof}

\KEntry{K03.14.uniform_comparison}{Uniform domination of tails}
\begin{theorem}
Suppose $g\ge0$ is locally properly integrable and $\int_A^\infty g$
converges. If $|f(x,t)|\le g(x)$ for $x\ge A$ and all $t\in T$,
and each $f$ section is locally properly integrable, then its improper
integrals converge uniformly on $T$, with
\[\sup_{t\in T}\left|\int_R^\infty f(x,t)\dd x\right|
\le\int_R^\infty g(x)\dd x\quad(R\ge A).
\]
\end{theorem}
\begin{proof}
On every regular $[R,S]$, the finite norm and order inequalities give
$|\int_R^Sf(x,t)\dd x|\le\int_R^Sg$ for every $t$.
The right-hand tails are uniformly small in $S\ge R$ by convergence
of $\int g$, so the uniform Cauchy criterion applies. Let $S\to\infty$
in the finite inequality to obtain the quantitative bound. Integrals
over the fixed initial interval exist separately by local integrability;
they need not obey a single domination bound to prove tail convergence.
Finite-endpoint versions have exactly the same proof on the corresponding
truncations.
\end{proof}

\KEntry{K03.14.uniform_dirichlet}{Uniform Dirichlet test}
\begin{theorem}
Let $f(\cdot,t)$ be locally Riemann integrable and
$A(x,t)=\int_a^xf(u,t)\dd u$. Suppose $|A(x,t)|\le C$ for all
$x\ge a$, $t\in T$. Let $g(\cdot,t)\ge0$ be $C^1$ and
nonincreasing, and suppose $\sup_{t\in T}g(R,t)\to0$.
Then $\int_a^\infty f(x,t)g(x,t)\dd x$ converges uniformly on $T$.
\end{theorem}
\begin{proof}
Products are locally integrable. The finite accumulation-parts proof
in \KRef{K03.10.dirichlet} works for each fixed $t$, with the same
constant: the primitive anchored at $R$ has modulus at most $2C$.
Thus
\[\left|\int_R^Sf(x,t)g(x,t)\dd x\right|
\le2C g(R,t)\le2C\sup_{u\in T}g(R,u).
\]
The right side tends to zero independently of $S\ge R$ and $t$.
The uniform Cauchy criterion proves the claim. Signed amplitudes of a
fixed sign are covered by replacing $g$ by its modulus when that modulus
has the same monotonicity and uniform limit property.
\end{proof}

\KEntry{K03.14.improper_continuity}{Continuity from uniform tails}
\begin{theorem}
If every proper truncation $F_R$ is continuous at $t_0$, and the
improper tails converge uniformly on a relative neighborhood $T_0$
of $t_0$, then $F(t)=\int_a^\infty f(x,t)\dd x$ is continuous at $t_0$.
\end{theorem}
\begin{proof}
Choose $R$ with $\sup_{t\in T_0}|F(t)-F_R(t)|<\eps/3$.
Continuity of this fixed $F_R$ gives
$|F_R(t)-F_R(t_0)|<\eps/3$ for all sufficiently close $t\in T_0$.
The triangle inequality gives
\[|F(t)-F(t_0)|\le|F(t)-F_R(t)|+
|F_R(t)-F_R(t_0)|+|F_R(t_0)-F(t_0)|<\eps.
\]
Joint continuity on each regular compact rectangle is one sufficient
condition for the truncation continuity, by the proper theorem.
\end{proof}

\KEntry{K03.14.improper_derivative}{Differentiation from uniform derivative tails}
\begin{theorem}
Let $[c,d]$ be a compact parameter interval. Suppose every truncation
$F_R$ is continuous on $[c,d]$ and differentiable inside, and proper
differentiation is valid there:
\[F_R'(t)=\int_a^R f_t(x,t)\dd x.
\]
Assume these derivative integrals converge uniformly on $[c,d]$ to
$G(t)$ (using endpoint extensions where needed), and $F_R(t_*)$ has a
finite limit at one $t_*\in[c,d]$. Then $F_R$ converges uniformly on
$[c,d]$ to an improper integral family $F$, differentiable inside with
$F'=G$. Joint continuity of $f,f_t$ on every regular compact rectangle
is a sufficient condition for the proper differentiation assumptions.
\end{theorem}
\begin{proof}
For any two late truncations, the norm mean-value inequality gives
\[\sup_{t\in[c,d]}|F_R(t)-F_S(t)|
\le|F_R(t_*)-F_S(t_*)|+(d-c)\sup_{u\in[c,d]}|F_R'(u)-F_S'(u)|.
\]
Both terms tend to zero, so completeness pointwise and the same uniform
Cauchy argument as above yield a uniform limit $F$. In particular all
original improper integrals now exist, not just the one at $t_*$.
Fix an interior $t$ and a nonzero $h$ with $t+h\in[c,d]$.
Apply the norm mean-value inequality to $F_S-F_R$ between these points
and let $S\to\infty$. Uniform convergence of derivatives gives
\[\left|\frac{F(t+h)-F(t)}h-
\frac{F_R(t+h)-F_R(t)}h\right|
\le\sup_{u\in[c,d]}|G(u)-F_R'(u)|.
\]
Hence the difference of the first quotient from $G(t)$ is at most this
uniform error plus
$|(F_R(t+h)-F_R(t))/h-F_R'(t)|+|F_R'(t)-G(t)|$.
Choose $R$ first to make the first and third terms small, then let
$h\to0$ for that fixed differentiable $F_R$. This proves $F'=G$.
It is the derivative-limit argument of \KRef{K03.12.differentiation}
spelled out for a continuous truncation parameter; uniform convergence
of the original integrands alone would not justify this conclusion.
\end{proof}

\KEntry{K03.14.parameter_order}{Integrating an improper family in the parameter}
\begin{theorem}
Suppose $f$ is jointly continuous on $[a,R]\times[c,d]$ for every
finite $R>a$, and $\int_a^\infty f(x,t)\dd x$ converges uniformly
for $t\in[c,d]$. Then
\[\int_c^d\left(\int_a^\infty f(x,t)\dd x\right)\dd t
=\int_a^\infty\left(\int_c^df(x,t)\dd t\right)\dd x,
\]
with a proper outer integral on the left and a convergent improper
outer integral on the right.
\end{theorem}
\begin{proof}
The truncations $F_R$ are continuous and converge uniformly, so $F$ is
continuous and properly integrable on $[c,d]$ by
\KRef{K03.08.uniform_stability}. Let
$H(x)=\int_c^df(x,t)\dd t$. Proper parameter continuity shows $H$
is continuous on each regular $[a,R]$, hence locally properly integrable.
For every $R$, the proper order theorem gives
$\int_a^RH(x)\dd x=\int_c^dF_R(t)\dd t$.
The right side tends to $\int_c^dF$ with error at most
$(d-c)\|F_R-F\|_\infty$. Thus the left partial integrals have a
finite limit, proving existence of the improper right side and the
claimed equality simultaneously. More generally the same proof works
if truncation order interchange and integrability are established by
other explicit proper hypotheses.
\end{proof}

\KEntry{K03.14.local_uniform}{Local uniformity is sufficient}
\begin{proposition}
On an open parameter interval $J$, uniform tail convergence on each
compact parameter interval implies continuity whenever truncations are
continuous. Uniform convergence of the derivative-integral tails on
each compact interval, with convergence of the original family at one
point of each such interval, gives differentiability there. If $J$ is
an interval and convergence is known at just one point, the base-point
condition propagates to all compact subintervals by overlapping intervals.
\end{proposition}
\begin{proof}
Every $t_0\in J$ is in the interior of some $[c,d]\subset J$.
Apply the preceding continuity or differentiation theorem on that
compact neighborhood. On overlaps the limiting functions agree because
they are limits of the same truncations; so do their derivative formulas.
For propagation from a fixed $t_*$ to any compact $K\subset J$, the
closed interval spanning $K\cup\{t_*\}$ lies in $J$. Uniform derivative
tails there and the one base point give convergence throughout by the
preceding theorem. No single bound uniform over an unbounded $J$ is
required for these local conclusions.
\end{proof}

\KEntry{K03.14.failure}{Pointwise convergence without parameter continuity}
\begin{example}
For $t\ge0$, $x\ge0$, put $f(x,t)=tE(-tx)$, so $f(x,0)=0$.
For each fixed $x$ this is continuous at zero (indeed jointly continuous
on finite rectangles). For $t>0$,
\[F_R(t)=\int_0^RtE(-tx)\dd x=1-E(-tR),\qquad
F(t)=1.
\]
At $t=0$ every partial integral and the improper value are zero, so $F$
is discontinuous. The tail for $t>0$ equals $E(-tR)$, whose supremum
over $0<t\le\delta$ is one for every $R$. Thus uniform convergence
fails on every neighborhood of zero, exactly where the continuity
theorem requires it. The finite computations use the proved exponential
derivative and the infinite computation uses $E(-tR)\to0$ for fixed $t>0$.
\end{example}

\KEntry{K03.14.boundary_terms}{A moving boundary supplies the derivative}
\begin{example}
Let $f(x,t)=1$ and $F(t)=\int_0^t1\dd x=t$, using orientation for
$t<0$. Thus $F'(t)=1$. The integrand parameter derivative is zero;
Leibniz's rule supplies the value entirely from the upper boundary term
$f(t,t)\,\frac{\dd t}{\dd t}=1$. Omitting that term would give the
incorrect derivative zero, even for this continuous constant integrand.
\end{example}
